% BANKS-SOURCE: pdf=75 print=65 kind=section id=6.2
\subsection{Heavy-fermion production in electron--positron annihilation}

% BANKS-SOURCE: pdf=75 print=65 kind=prose
In this section we will calculate the leading-order contribution, in massive QED, to the
annihilation cross section for a light charged particle (electrons and positrons) into a
heavier one. We will call the heavier particles muons, but the calculation is valid for
$\tau$ leptons as well. As a consequence of asymptotic freedom, it can also be used for
calculating the $e^+e^-\to$ hadrons cross section above the QCD confinement scale (see
Chapter 8 for a definition). In that case, the rigorous use of the calculation involves an
analytic continuation from Euclidean space~\cite{banks-ref-32}, which we will not delve into.
Since we work with massive photons, our calculation can also be adapted to calculate the part
of the full cross section that comes from $Z$-boson, rather than photon, exchange. We
need only substitute the correct couplings to the $Z$ boson, and set $\mu$ equal to $m_Z$ in the
calculation below. We will learn how to calculate the $Z$ couplings in Chapter 8. Given
the format of this book, we cannot possibly do justice to the full extent of the physics
of this process. The reader is urged to consult the wonderful section in Chapter 5 of
Peskin and Schroeder~\cite{banks-ref-33} for a detailed description of it.

% BANKS-SOURCE: pdf=76 print=66 kind=figure id=6.2
\begin{figure}[htbp]
  \centering
  % Native TikZ reconstruction of Figure 6.2.
\begin{tikzpicture}[x=1cm,y=1cm,>=Stealth,baseline=(current bounding box.center)]
  \tikzset{BanksFermion/.style={line width=.65pt,postaction={decorate},
    decoration={markings,mark=at position .5 with {\arrow{Stealth[length=5pt]}}}},
    BanksPhoton/.style={line width=.65pt,decorate,
      decoration={snake,amplitude=5pt,segment length=11pt}}}
  % incoming e^+ / e^- : one continuous fermion line, top-left to bottom-left
  \draw[BanksFermion] (-2.3,1.25) -- (-.85,0);
  \draw[BanksFermion] (-.85,0) -- (-2.3,-1.25);
  % massive photon
  \draw[BanksPhoton] (-.85,0) -- (.85,0);
  % outgoing mu^- / mu^+ : continuous fermion line, bottom-right to top-right
  \draw[BanksFermion] (2.3,-1.25) -- (.85,0);
  \draw[BanksFermion] (.85,0) -- (2.3,1.25);
  \node[above left=-2pt] at (-2.3,1.25) {$\ee^+$};
  \node[below left=-2pt] at (-2.3,-1.25) {$\ee^-$};
  \node[above right=-2pt] at (2.3,1.25) {$\mu^+$};
  \node[below right=-2pt] at (2.3,-1.25) {$\mu^-$};
\end{tikzpicture}

  \caption{Heavy-fermion production in $e^+,e^-$ annihilation.}
  \label{fig:6.2}
\end{figure}

% BANKS-SOURCE: pdf=76 print=66 kind=prose
There is only one diagram (Figure~\ref{fig:6.2}) in leading order, and its value is

% BANKS-SOURCE: pdf=76 print=66 kind=equation id=6.6
\begin{equation}
  \mathcal{M}=(-\ii e)^2\bar u(k_-,r_- )\gamma^\kappa v(k_+,r_+ )
  \frac{-\ii(\eta_{\kappa\lambda}+q_\kappa q_\lambda/\mu^2)}
  {q^2+\mu^2-\ii\epsilon}
  \bar v(p_+,s_+)\gamma^\lambda u(p_-,s_-).
  \label{eq:6.6}
\end{equation}
Here $q=p_++p_-=k_++k_-$, and the labels indicate incoming and outgoing fermion
momenta and spins in what I hope is an obvious manner. The labels are related to the
scattering energy and angle in the center-of-mass frame by Figure~\ref{fig:6.3}.

We note the identity
\begin{equation*}
  \bar v(p_+,s_+)(p_++p_-)_{\lambda}\gamma^\lambda u(p_-,s_-)
  =(-m_e+m_e)\bar v(p_+,s_+)u(p_-,s_-)=0,
\end{equation*}
which eliminates the longitudinal part of the gauge-boson propagator from the
amplitude.

We will now calculate the spin-averaged total cross section for this process where we
sum over initial and average over final spins. This corresponds to the simplest experi-
mental situation, in which the initial beams are unpolarized and no spin measurements

% BANKS-SOURCE: pdf=76 print=66 kind=figure id=6.3
\begin{figure}[htbp]
  \centering
  % Native TikZ reconstruction of Figure 6.3.
\begin{tikzpicture}[x=1cm,y=1cm,>=Stealth,baseline=(current bounding box.center)]
  \tikzset{BanksFermion/.style={line width=.62pt,postaction={decorate},
    decoration={markings,mark=at position .5 with {\arrow{Stealth[length=5pt]}}}}}
  % ---- momentum labels for the outgoing muons -------------------------
  \node at (0,2.21)
    {$K_+^\nu=\left(E_{\mathbf k},0,-|\mathbf k|\sin\theta,-|\mathbf k|\cos\theta\right)$};
  \node at (0,-2.32)
    {$K_-^\nu=\left(E_{\mathbf k},0,|\mathbf k|\sin\theta,|\mathbf k|\cos\theta\right)$};
  % ---- incoming e^- (from the left) and e^+ (from the right) ----------
  \draw[BanksFermion] (-3.16,0) -- (0,0);
  \draw[BanksFermion] (3.12,0) -- (0,0);
  % ---- outgoing muons along a single straight line through the vertex --
  \draw[BanksFermion] (0,0) -- (2.39,1.45);
  \draw[BanksFermion] (0,0) -- (-2.40,-1.45);
  \fill (0,0) circle (2.2pt);
  % ---- angle tick marks ------------------------------------------------
  \draw[line width=.62pt] (1.10,-0.03) -- (1.10,0.11);
  \draw[line width=.62pt]
    ($(31.2:1.10)+(121.2:-0.04)$) -- ($(31.2:1.10)+(121.2:0.10)$);
  \draw[line width=.62pt] (-1.10,0.03) -- (-1.10,-0.11);
  \draw[line width=.62pt]
    ($(211.2:1.10)+(301.2:-0.04)$) -- ($(211.2:1.10)+(301.2:0.10)$);
  \node at (1.02,0.30) {$\theta$};
  \node at (-1.11,-0.31) {$\theta$};
  % ---- external particle labels ---------------------------------------
  \node[left] at (-3.24,0) {$\ee^-$};
  \node[right] at (3.20,0) {$\ee^+$};
  \node[above right=-2pt] at (2.39,1.45) {$\mu^+$};
  \node[below left=-2pt] at (-2.40,-1.45) {$\mu^-$};
  \node at (-3.10,0.62) {$p_-^\nu=(E_{\mathbf p},0,0,-|\mathbf p|)$};
  \node at (3.22,0.62) {$p_+^\nu=(E_{\mathbf p},0,0,|\mathbf p|)$};
\end{tikzpicture}

  \caption{$\mu^+,\mu^-$ pair production by annihilation of massless $e^+,e^-$ in the center-of-mass frame.}
  \label{fig:6.3}
\end{figure}

% BANKS-SOURCE: pdf=77 print=67 kind=prose
are made on the final particles. Polarized beams and polarization detectors are power-
ful experimental tools in $e^+e^-$ annihilation. The reader is urged again to consult Peskin
and Schroeder~\cite{banks-ref-33} for a detailed description of polarization amplitudes. At this point
the reader should also turn to the appendix on Diracology. He/she is urged to take all
of the results stated there as exercises, even those whose proofs are given. Mastery of
this technology is an important part of the repertoire of any respectable field theorist.

\BanksImplicitHook{I-CH06-002}

Using the polarization sum identities of Diracology, the spin-averaged squared
amplitude is

% BANKS-SOURCE: pdf=77 print=67 kind=equation id=6.7
\begin{equation}
  \begin{aligned}
  \frac14\sum_{r_\pm,s_\pm}|\mathcal{M}|^2
  ={}&\frac{e^4}{4\lvert(q^2+\mu^2-\ii\epsilon)\rvert^2}
  \operatorname{tr}\!\left[(\slashed p_+-m_e)\gamma^\lambda
  (\slashed p_-+m_e)\gamma^\kappa\right]\\[-2pt]
  &\quad\times\operatorname{tr}\!\left[(\slashed k_-+m_\mu)\gamma_\lambda
  (\slashed k_+-m_\mu)\gamma_\kappa\right].
  \end{aligned}
  \label{eq:6.7}
\end{equation}

Now we use the trace formulae of the appendix to write this as

% BANKS-SOURCE: pdf=77 print=67 kind=equation id=6.8
\begin{equation}
  \frac14\sum_{r_\pm,s_\pm}|\mathcal{M}|^2
  =\frac{e^4}{4\lvert(q^2+\mu^2-\ii\epsilon)\rvert^2}
  \left[(p_-\!\cdot k_-)(p_+\!\cdot k_+)+(p_-\!\cdot k_+)(p_+\!\cdot k_-)
  -m_\mu^2(p_-\!\cdot p_+)\right].
  \label{eq:6.8}
\end{equation}
We have dropped a term proportional to $m_e^2$ (which the reader should calculate), since
these formulae are always used at center-of-mass energy above twice the muon mass.

\BanksImplicitHook{I-CH06-001}

We evaluate the scattering cross section in the center-of-mass frame (which is not
quite the lab frame in an $e^+e^-$ collider, because of bremsstrahlung). There, according
to Figure~\ref{fig:6.3},
\begin{equation*}
  (p_++p_-)^2=-4E^2;\qquad (p_+\!\cdot p_-)=-2E^2;
\end{equation*}
\begin{equation*}
  (p_-\!\cdot k_-)=(p_+\!\cdot k_+)
  =-E^2-E\sqrt{E^2-m_\mu^2}\cos\theta;
\end{equation*}
\begin{equation*}
  (p_-\!\cdot k_+)=(p_+\!\cdot k_-)
  =-E^2+E\sqrt{E^2-m_\mu^2}\cos\theta.
\end{equation*}

We have again made the approximation that the electrons are massless and travel at the
speed of light. $E$ is the energy and absolute value of the momentum of the electron or
positron. The square root in the above formulae is the absolute value of the muon or
anti-muon momentum.

We now consult the appendix on cross sections. The following angular formulas use
the $D=4$ specialization, with three spatial momentum components. Using the
two-to-two formula with
$E_A=E_B=E$, and $|v_A-v_B|=2$ to write

% BANKS-SOURCE: pdf=77 print=67 kind=equation id=6.9
\begin{equation}
  \frac{d\sigma}{d\Omega}=\frac{\sqrt{1-m_\mu^2/E^2}}{1024\pi^2}
  \sum|\mathcal{M}|^2.
  \label{eq:6.9}
\end{equation}

% BANKS-SOURCE: pdf=77 print=67 kind=equation id=6.10
\begin{equation}
  \frac{d\sigma}{d\Omega}
  =\frac{\alpha^2}{(4E^2-\mu^2)^2}\sqrt{1-\frac{m_\mu^2}{E^2}}
  \left[(E^2+m_\mu^2)+(E^2-m_\mu^2)\cos^2\theta\right].
  \label{eq:6.10}
\end{equation}
We have introduced the fine-structure constant $\alpha\equiv e^2/4\pi\approx1/137$.

% BANKS-SOURCE: pdf=78 print=68 kind=prose
There are three interesting things to note about this cross section. The first is the
square-root turn on of the cross section near threshold and the second is the angular
distribution of the reaction products. Both the square root and the angular distribu-
tion are characteristic of spin-$\frac12$ particles. Observations of thresholds like this in $e^+e^-$
annihilation signal new quarks and leptons and give us one of the pieces of evidence
that they all have spin $\frac12$. If we are ever lucky enough to see a supersymmetric partner
of a quark or lepton with spin zero, we will see something quite different. The electro-
magnetic current for these particles carries a factor of momentum, so the turn on of the
cross section above threshold is slower by an extra square root. The angular distribution
is also different. You should do Problem 6.4, to appreciate these differences.

Finally, if we calculate the total cross section

% BANKS-SOURCE: pdf=78 print=68 kind=equation id=6.11
\begin{equation}
  \sigma_{\rm tot}=\frac{16\alpha^2\pi}{3(4E^2-\mu^2)^2}
  \sqrt{1-\frac{m_\mu^2}{E^2}}\left(E^2+\frac12m_\mu^2\right).
  \label{eq:6.11}
\end{equation}
and take the limit $E\gg$ all mass scales, then we get the scale-invariant result

% BANKS-SOURCE: pdf=78 print=68 kind=equation id=6.12
\begin{equation}
  \sigma_{\rm tot}=\frac{\pi\alpha^2}{3E^2}.
  \label{eq:6.12}
\end{equation}
The fact that this is scale-invariant seems obvious, but it is a clue to the very essence
of what a quantum field theory is. You’ll have to hold your breath till the chapter on
renormalization to find out.

\BanksImplicitHook{I-CH06-003}
